\subsection{Punto de partida: h(n)=0}

El archivo inicial trae \texttt{cornersHeuristic} devolviendo siempre \texttt{0}. Como senala la
guia, esta funcion es admisible por definicion (nunca sobreestima nada) pero no aporta ninguna
informacion util para orientar la busqueda: con ella, A* se comporta exactamente igual que UCS.

\subsection{Heuristica basica}

Se implemento primero la formula que sugiere la guia como primera aproximacion: la distancia
Manhattan a la esquina pendiente mas lejana.

$$h(n) = \max_{c \in C_p} d_M(n, c)$$

donde $C_p$ es el conjunto de esquinas que todavia no se han visitado en el estado $n$ (si ya no
queda ninguna esquina pendiente, $h(n) = 0$).

\begin{verbatim}
def cornersHeuristicBasica(state, problem):
    position, visited = state
    corners = problem.corners
    pendientes = [c for c, v in zip(corners, visited) if not v]
    if not pendientes:
        return 0
    distancias = (abs(position[0]-c[0]) + abs(position[1]-c[1])
                  for c in pendientes)
    return max(distancias)
\end{verbatim}

\subsection{Heuristica propuesta}

Siguiendo la sugerencia de la guia de analizar ``si es posible disenar una heuristica mas
informativa'', se generalizo la heuristica basica con el mismo patron de diametro que ya se habia
usado en \texttt{foodHeuristic} (Actividad 11): en vez de mirar solo la distancia de Pac-Man a la
esquina pendiente mas lejana, se considera el diametro --la mayor distancia Manhattan entre
cualquier par-- del conjunto formado por la posicion actual mas todas las esquinas pendientes.

$$h(n) = \max\Big(\max_{c \in C_p} d_M(n, c),\ \max_{c_i, c_j \in C_p} d_M(c_i, c_j)\Big)$$

Es mas informativa que la basica porque tambien tiene en cuenta que, para visitar dos esquinas
pendientes lejanas entre si, Pac-Man tiene que recorrer al menos la distancia real entre ellas en
algun momento del camino --no solo la distancia desde su posicion actual hasta la mas lejana de
las dos--.

\begin{verbatim}
def cornersHeuristic(state, problem):
    position, visited = state
    corners = problem.corners
    pendientes = [c for c, v in zip(corners, visited) if not v]
    if not pendientes:
        return 0
    puntos = [position] + pendientes
    maxDist = 0
    for i in range(len(puntos)):
        for j in range(i + 1, len(puntos)):
            p1, p2 = puntos[i], puntos[j]
            d = abs(p1[0]-p2[0]) + abs(p1[1]-p2[1])
            if d > maxDist:
                maxDist = d
    return maxDist
\end{verbatim}

A diferencia de \texttt{foodHeuristic}, aqui no se uso \texttt{problem.heuristicInfo} como cache:
como mucho hay 4 esquinas (nunca mas de 6 pares), asi que el costo de recalcular todas las
distancias entre pares en cada llamada es insignificante --muy distinto al caso de
\texttt{FoodSearchProblem}, donde puede haber decenas de alimentos y cientos de pares--. Cachear
aqui anadiria la misma sobrecarga de diccionario que en la Actividad 11 no compenso, sin ningun
beneficio real dado lo pequeno del conjunto.

\subsection{Admisibilidad}

\textbf{Heuristica basica.} El costo real para visitar todas las esquinas pendientes es, como
minimo, el costo de llegar hasta la esquina pendiente mas lejana --en algun momento del recorrido
Pac-Man tiene que llegar hasta ella--, y la distancia Manhattan nunca sobreestima una distancia
real en una cuadricula con paredes (cada paso real cambia $x$ o $y$ en una unidad; la distancia
Manhattan es el minimo numero de esos pasos sin contar paredes). Por lo tanto $h(n) \le h^*(n)$.

\textbf{Heuristica propuesta.} El recorrido que visita, en cualquier orden, la posicion $n$ y
todas las esquinas de $C_p$ tiene una longitud real que es, como minimo, la distancia real entre
el par de puntos mas separado entre si de $\{n\} \cup C_p$ --en algun momento del recorrido
Pac-Man pasa por ambos puntos de ese par--. Como la distancia Manhattan nunca sobreestima la
distancia real, el diametro en Manhattan de $\{n\} \cup C_p$ es una cota inferior valida:
$h(n) \le h^*(n)$.

Ambas cotas se verificaron tambien de forma empirica en el estado inicial de \texttt{tinyCorners}
(ver el script de la Actividad 8 en \texttt{experimentos/}): $h_{basica}(inicio) = 6$ y
$h_{propuesta}(inicio) = 12$, ambos por debajo del costo optimo real conocido (22, obtenido con
UCS en la Actividad 7).

\subsection{Consistencia (no exigida, se incluye por completitud)}

Para cualquier sucesor $n'$ de $n$ (un paso ortogonal, costo 1), si el par de puntos que define
$h(n)$ sigue intacto en $n'$ (ninguno de los dos era la esquina que se acaba de visitar al pasar
de $n$ a $n'$), entonces por la desigualdad triangular de Manhattan
$d_M(n, p) \le d_M(n, n') + d_M(n', p) = 1 + d_M(n', p) \le 1 + h(n')$ para el punto $p$
correspondiente. Si uno de los dos puntos del par era justo la posicion $n$ (que cambia a $n'$) o
una esquina que se acaba de visitar, el argumento se reduce al caso trivial
$d_M(n, n') = 1 \le 1 + h(n')$, valido porque $h(n') \ge 0$ siempre. En consecuencia,
$h(n) \le c(n, n') + h(n')$ se cumple en todos los casos, y $h(\text{meta}) = 0$ (no quedan
esquinas pendientes). Ambas heuristicas son, entonces, tambien consistentes.

\subsection{Procedimiento y resultados}

Se ejecuto A* con las tres heuristicas sobre \texttt{tinyCorners} (ver el script de la
Actividad 8 en \texttt{experimentos/}).

\begin{table}[H]
\centering
\caption{A* sobre CornersProblem con distintas heuristicas (tinyCorners)}
\label{tab:act8-heuristicas}
\begin{tabular}{lrrrr}
\toprule
\textbf{Heuristica} & \textbf{Costo} & \textbf{Longitud} & \textbf{Nodos expandidos} & \textbf{Tiempo (s)} \\
\midrule
$h(n)=0$              & 22 & 22 & 295 & 0.001486 \\
Basica                & 22 & 22 & 147 & 0.000832 \\
Propuesta             & 22 & 22 & 119 & 0.000802 \\
\bottomrule
\end{tabular}
\end{table}

Las tres estrategias encuentran el mismo costo optimo (22, el mismo que dio UCS en la Actividad
7), confirmando que ninguna de las dos heuristicas disenadas sobreestima el costo real. La
heuristica basica reduce los nodos expandidos de 295 a 147 (49.8\,\% menos que $h=0$); la
heuristica propuesta reduce aun mas, a 119 (59.7\,\% menos que $h=0$, y 19.0\,\% menos que la
basica), confirmando que el termino adicional de distancia entre pares de esquinas pendientes si
aporta informacion util a la busqueda.

\subsection{Pregunta de análisis}

\textbf{Explique por qué su heurística para \texttt{CornersProblem} es admisible.}

Porque en ambos casos el valor que devuelve la heuristica es una distancia Manhattan (o el maximo
de varias distancias Manhattan) entre puntos que, de todas formas, Pac-Man tiene que recorrer en
algun momento para completar el objetivo --llegar a cada esquina pendiente, y en el caso de la
heuristica propuesta, tambien recorrer la distancia entre pares de esquinas pendientes entre si--.
Como la distancia Manhattan es, por construccion, una cota inferior de cualquier distancia real
en una cuadricula con obstaculos (un camino real nunca es mas corto que la distancia en linea
recta por cuadricula, porque cada paso real cambia como maximo una coordenada en una unidad), el
valor devuelto nunca puede ser mayor que el costo optimo real restante $h^*(n)$. No basta con decir
``la heuristica funciona''\ --lo que hace que sea admisible es, especificamente, que cada termino
del maximo es una distancia Manhattan entre dos puntos que forman parte, necesariamente, de
cualquier solucion real desde ese estado.
